\documentclass[11pt,a4paper]{article}

\usepackage[utf8]{inputenc}
\usepackage[T1]{fontenc}
\usepackage[magyar]{babel}

\usepackage[margin=2.2cm]{geometry}
\usepackage{xcolor}
\usepackage{titlesec}
\usepackage{tcolorbox}
\tcbuselibrary{skins, breakable}
\usepackage{enumitem}
\usepackage{booktabs}
\usepackage{graphicx}
\usepackage{hyperref}
\usepackage{tabularx}
\usepackage{amsmath}
\usepackage{microtype}
\usepackage{lmodern}

% Colors
\definecolor{primary}{RGB}{30, 80, 150}
\definecolor{accent}{RGB}{40, 140, 90}
\definecolor{danger}{RGB}{180, 50, 50}

\hypersetup{
    colorlinks=true,
    linkcolor=primary,
    urlcolor=primary,
    pdftitle={SUMO RL Traffic Light Control — Methodology \& Research Findings}
}

% Section formatting
\titleformat{\section}{\Large\bfseries\color{primary}}{\thesection.}{1em}{}[\titlerule]
\titleformat{\subsection}{\large\bfseries\color{accent}}{\thesubsection}{1em}{}

% Custom boxes
\newtcolorbox{infobox}[1][]{
  colback=primary!5!white,
  colframe=primary,
  fonttitle=\bfseries,
  title=#1,
  boxrule=0.8pt,
  arc=3pt,
  breakable,
  top=2mm, bottom=2mm
}

\newtcolorbox{warningbox}[1][]{
  colback=danger!5!white,
  colframe=danger,
  fonttitle=\bfseries,
  title=#1,
  boxrule=0.8pt,
  arc=3pt,
  breakable,
  top=2mm, bottom=2mm
}

\newtcolorbox{codebox}[1][]{
  colback=gray!5!white,
  colframe=gray!60!black,
  fonttitle=\bfseries\ttfamily,
  title=#1,
  boxrule=0.5pt,
  arc=2pt,
  breakable
}

\newcommand{\uv}[1]{\glqq#1\grqq}
\newcommand{\code}[1]{\texttt{\color{black!80}#1}}
\setlength{\parindent}{0pt}
\setlength{\parskip}{0.7em}

\begin{document}

\begin{center}
    \vspace*{1em}
    {\Huge\bfseries\color{primary} SUMO RL Traffic Light Control}\vspace{0.5em}\\
    {\Large\bfseries Methodology \& Research Findings}\vspace{0.8em}\\
    {\large Comprehensive Documentation for Paper Writing}
    \vspace{1.5em}
\end{center}

\begin{infobox}[Meta adatok]
\begin{itemize}[leftmargin=*, nosep]
    \item \textbf{Generálva:} 2026-03-14
    \item \textbf{Adatforrás:} \code{metric\_pca\_test\_v2/} (527,747 sample)
    \item \textbf{Szimuláció:} SUMO (Simulation of Urban MObility) + traci/libsumo
    \item \textbf{RL framework:} Stable-Baselines3 QR-DQN (Quantile Regression DQN)
    \item \textbf{Hálózat:} \code{mega\_catalogue\_v2.net.xml} — 21 jelzőlámpás csomópont (4×5 + 2 osztott)
\end{itemize}
\end{infobox}

\textbf{EZ A DOKUMENTUM CÉLJA:} Mindent tartalmaz, amit a cikk megírásához tudni kell:
Mi a rendszer és hogyan működik, honnan jönnek a mérőszámok, miért pont ezeket a metrikákat választottuk a reward-hoz, a normalizáció, háló architektúra, hiperparaméter-keresés, új baseline stratégiák, a training pipeline, valamint a curriculum learning megfontolásai.

\tableofcontents
\newpage

\section{System Architecture — A Rendszer Felépítése}

\textbf{A Nagy Kép:} Van egy városi úthálózatunk (SUMO szimulátorban), 21 jelzőlámpás csomóponttal. Mindegyik csomóponthoz tartozik egy \uv{ágens} — egy mesterséges intelligencia, amely megtanulja, mikor melyik fázist (zöld irányt) kapcsolja be. Az ágensek egymástól függetlenül tanulnak (nincs kommunikáció köztük), de ugyanazt a hálózati architektúrát használják, külön-külön súlyokkal.

\subsection{Multi-Agent RL Setup}
\begin{itemize}
    \item 21 csomópont, mindegyik önálló RL ágens. Csomópont azonosítók: \code{R1C1\_C ... R4C5\_C} (4×5 rács), \code{R4C2\_CE/R4C2\_CW} (kettéosztott).
    \item Algoritmus: QR-DQN (Quantile Regression Deep Q-Network).
\end{itemize}

\begin{infobox}[Miért QR-DQN és nem sima DQN?]
A hagyományos DQN egyetlen számot (átlagos várható érték) becsül minden cselekvésre. A QR-DQN viszont az EGÉSZ ELOSZLÁST becsli — tehát tudja, hogy \uv{ez a cselekvés átlagosan jó, de nagy a variancia}. Ez fontos, mert a forgalom természeténél fogva sztochasztikus: ugyanaz a jelzésprogramozás néha jól működik, néha nem.
\end{infobox}

\textbf{Miért független tanulás (nem kooperatív)?}
Kooperatív tanulás (ahol az ágensek kommunikálnak) exponenciálisan növekvő állapotteret jelent 21 ágensre. A független tanulás skálázhatóbb, és a gyakorlatban jól működik, mert a szomszédos csomópontok hatása az observation-ön keresztül implicit megjelenik.

\subsection{Observation Space — Mit Lát Az Ágens?}
Minden lépésben az ágens a következő információkat kapja:
\begin{itemize}
    \item \textbf{phase} (1 szám): Melyik zöld fázis van éppen bekapcsolva (0, 1, 2, \dots).
    \item \textbf{phase\_timer} (1 szám, 0.0 és 1.0 között): Mennyi ideje tart az aktuális fázis, normalizálva (pl. 120 lépés = 600 sec = 1.0).
    \item \textbf{occupancy} (N szám): Detektorok foglaltsága 0 és 1 között (0=üres, 1=foglalt). A legfontosabb bemenet.
    \item \textbf{flow} (N szám): Hány jármű haladt át a detektoron a lépés alatt.
\end{itemize}

\subsection{Action Space — Mit Csinálhat Az Ágens?}
\code{Discrete(num\_phases)} — annyi lehetséges cselekvés, ahány zöld fázis van. (Pl. 4 fázis esetén \{0, 1, 2, 3\}). Ezenkívül érvényesül az \textbf{Intergreen rendszer} (átmeneti mátrix beiktatása piros és sárga időkkel fázisváltások között) és a \textbf{Minimum Zöld Idő} biztosítása (pl. legalább 5 mp).

\subsection{Step Mechanics}
\code{delta\_time = 5 SUMO sec / RL lépés}. Az ágens választ egy akciót, ezt követi 5 db 1 másodperces SUMO szimulációs substep. Ezek alatt a metrikákat összegezzük, majd az akkumulált értékek átlagát képezzük.

\section{Traffic Generation — Hogyan Generáljuk A Forgalmat?}

\subsection{Lane-Level Random Traffic \& Target+Spread Randomizálás}
Nem a SUMO beépített \code{randomTrips.py}-t használjuk. Helyette sáv-szintű véletlenszerű forgalmat generálunk.
Az új megoldás epizód-szinten két fázisban randomizál a \uv{nagy számok törvénye} elkerülésére:
\begin{itemize}
    \item \textbf{Target:} Az epizód összforgalmának nagyságrendje (pl. 100, 500, 900 veh/h).
    \item \textbf{Spread:} Az irányok közötti aszimmetria mértéke.
\end{itemize}
Ezzel az ágens kifejezetten egyenletes, de aszimmetrikus; könnyű és nehéz szituációkat egyaránt megtapasztal.

\subsection{Episode Duration Rotation \& Warm-Up kikapcsolása}
Minden epizódban véletlenszerűen kiválasztunk egy időtartamot: 15 perctől 2 óráig. Ez meggátolja a temporális overfitting-et. Ezen kívül a random \textit{Warm-up} funkció ki lett kapcsolva, mert a megörökölt kezdeti \uv{mesterséges} dugó nehezítette a tanulást a \uv{memória nélküli} ágens számára.

\section{Metric Collection \& SUMO API}
\subsection{Raw SUMO Metrics}
\begin{itemize}
    \item \textbf{TotalWaitingTime [sec/step]:} Sávon várakozó járművek összes ideje. \textbf{Fő reward metrika.}
    \item \textbf{TotalCO2 [mg/s/step]:} Sáv összes járművének pillanatnyi CO2 kibocsátása (HBEFA modell). \textbf{Reward része.}
    \item \textbf{AvgCO2:} Konstantishoz közeli, HIBÁS tanulási szignál.
\end{itemize}

\begin{warningbox}[A TravelTime Csapdája (Sentinel Values)]
A SUMO \code{getTraveltime()} üres/bedugult sáv esetén $\approx$ 200,000 és 5,000,000 közötti sentinel értékeket ad vissza. A PCA ezt izolálja, mint ZAJDIMENZIÓ. Szűrés nélkül használhatatlan.
\end{warningbox}

\section{PCA Analysis — Változók Kiválasztása}
A PCA elemzés szerint:
\begin{itemize}
    \item \textbf{PC1 — Congestion (44.8\%):} \code{TotalWaitingTime}, \code{QueueLength}, stb. erősen korreláltak (redundánsak). 
    \item \textbf{PC2 — Travel Time (22.6\%):} Független torlódástól, de sentinel torzítást tartalmaz.
\end{itemize}
\textbf{Eredmény:} \code{TotalWaitingTime} (PC1 dom.) és \code{TotalCO2} használata az optimális a jutalom-függvény számára.

\section{Normalization \& Reward Function}
\subsection{Log-Sigmoid Normalizáció}
Az online data streaming miatt a \textbf{Log-Sigmoid} normalizációs stratégia bizonyult a leghatékonyabbnak (quality=1.354). Két konstansra támaszkodik metrikánként ($\mu, \sigma$).
Formula: $score = 1 / (1 + \exp(-((\log(x + \varepsilon) - \mu) / \sigma)))$.

\textbf{Reward Kombináció:} Az ágens a két metrika inverz-szigmoidálisan leképezett értékének átlagát kapja. Üres csomópontra szigorúan $0.0$ jutalom van garantálva a hamis pozitív megerősítés elkerülésére.

\section{Neural Network, Sweep \& Backends}
\subsection{QR-DQN és Sweep}
A háló \code{MultiInputPolicy} alapú, amely a W\&B sweep optimalizáció során hangolható paraméterekkel (\code{num\_layers}, \code{layer\_size}, tanulási ráta stb.) alakul ki. Mentés: \code{.zip} és \code{ONNX} formátumban.
\subsection{Backend (LIBSUMO vs TRACI)}
\begin{itemize}
    \item \textbf{libtraci / libsumo:} Natív kliens, szimpla futás révén (30-50\%) gyorsabb, azonban \code{libsumo} esetenként singleton korlátokkal bír. 
    \item \textbf{Traci:} Párhuzamosítható és elengedhetetlen a GUI-hoz / evaluációs szkriptekhez.
\end{itemize}
\textbf{SUMO Subscriptions:} A batch lekérdezések (pl. detektor adatok) implementációja mindkét backend-en drámaian meggyorsítja az adatgyűjtést.

\section{Baseline Evaluation — Összehasonlító Stratégiák}
Négy (baseline) stratégia került implementálásra a kutatás során az összehasonlításhoz (\code{baseline\_actuated.py}):
\begin{enumerate}
    \item \textbf{Static:} Beégetett \code{.net.xml} alapú fix ciklus, adaptáció nélkül.
    \item \textbf{Actuated:} SUMO beépített aktuált / szenzorvezérelt fix-sorrendű szimulációja (maxGap = 3.1s).
    \item \textbf{Random:} Determinisztikus de random fázist választó döntéshozó. Worst-case scenario.
    \item \textbf{Max-Pressure:} A közlekedésmérnöki irodalom jelenlegi standardja. Minden döntési ciklusban kiszámítja a \uv{nyomást} (a kapott zöldágakon megállt járművek összességét), és a legmagasabb értékűt részesíti előnyben.
\end{enumerate}
Ezek a stratégiák (\code{random} és \code{max-pressure}) teljesen ugyanazt az értékelési rendszert és RL környezetet (transition logic, reward, env) kapják meg a fair összehasonlítás érdekéhez. W\&B fiókkal könnyen követhetőek egymás mellett a grafikonokon.

\section{Training Pipeline \& Curriculum Learning}
\subsection{Kivitelezés és CLI argumentumok}
Két módon futtatható a rendszer: GUI (Tkinter), amellyel subprocessként (\code{main\_headless.py}) aktiváljuk, vagy natívan shellből/Dockerből paraméterekkel (\code{--learning-rate}, \code{--num-layers}).
\textbf{Prioritás:} Sweep config > CLI args > YAML defaults.

\subsection{Curriculum Learning megfontolásai (Domain Randomization)}
Fokozatos nehézség (forgalom) növelés komoly \textbf{Catastrophic Forgetting} hátulütőket eredményez a Limitált DQN replay buffer mérete miatt, emellett "Exploration Gap" jelenséget is produkálhat, ahol az ágens nem ismeri meg a nehéz stratégiákat a kezdeti kapott sok zöld jelzés miatt. 
Helyette a \textbf{Domain Randomization (Target+Spread)} nyújt stabilabb végeredményt aszimmetrikus és véletlengenerált könnyű/nehéz epizódjaival. Nincs szükség külső \code{flow\_min/max} módosításra, a diverzitás az epizód szintű logika miatt automatikus.

\section{Key Conclusions \& Bibliography}
\subsection{Összegzés}
\begin{itemize}
    \item \textbf{PCA Selection:} A 12 indikátorból 3 zajmentes komponens emelkedik ki: \code{WaitingTime} \& \code{TotalCO2}.
    \item \textbf{Normalization:} Az online stream környezethez hangolt pre-kalibrált Log-Sigmoid lett beépítve.
    \item \textbf{Baseline Eval:} Létrejött az objektív és ipari \code{max-pressure} és actuated baseline összehasonlító analitika a W\&B-en belül is.
    \item \textbf{Traffic Generation:} Lane-level target+spread logika nyújt hatékony generalizációt "curriculum learning" helyett.
\end{itemize}
\subsection{Kulcs Hivatkozások}
\begin{itemize}
    \item \textit{SUMO} (Lopez et al., 2018), \textit{QR-DQN} (Dabney et al., 2018), \textit{HBEFA 4.2}.
    \item \textit{Max-Pressure Control} (Varaiya 2013; Wei et al. 2019)
    \item \textit{Stable-Baselines3} (Raffin et al., 2021)
\end{itemize}

\end{document}
